\section{Full-panel results}\label{app:fullpanel}

The primary sample is truncated at the backfill horizon (\noaaEnds). The same scans
run to the data maximum (\bWeeksFull{} weeks) are reported here for completeness.
They were the runs first recorded in the decision log, and they change no
conclusion.

\begin{table}[h]
\centering\small
\caption{Model ladder, full panel.}
\begin{tabular}{llrrrrr}
\toprule
$h$ (wk) & model & $n$ & MAE & RMSE & 80\% cov. & skill vs M0 \\
\midrule
  1 & M0 & 288 & 1.997 & 4.371 & 82.3\% & -- \\
  1 & M1 & 288 & 2.023 & 4.473 & 82.3\% & $-$1.3\% \\
  1 & M2 & 288 & 2.205 & 4.542 & 83.0\% & $-$10.4\% \\
  4 & M0 & 282 & 3.947 & 7.090 & 73.8\% & -- \\
  4 & M1 & 282 & 4.168 & 7.444 & 72.0\% & $-$5.6\% \\
  4 & M2 & 282 & 4.705 & 7.378 & 69.1\% & $-$19.2\% \\
\bottomrule
\end{tabular}

\end{table}

\begin{table}[h]
\centering\small
\caption{FWL partial effects, full panel. The largest $|t|$ here (\bMaxTFull)
belongs to the unresolved-destination band and falls to $+0.24$ on the primary
sample. It was produced by the contaminated tail.}
\begin{tabular}{llrrrrr}
\toprule
 & & \multicolumn{2}{c}{$h=1$} & \multicolumn{2}{c}{$h=4$} & \\
\cmidrule(lr){3-4} \cmidrule(lr){5-6}
signal & prior & HAC $t$ & partial $R^2$ & HAC $t$ & partial $R^2$ & year cons. \\
\midrule
  gas in transit, EU-bound & $+$ & +0.08 & 0.000 & +1.03 & 0.004 & 75\% \\
  gas in transit, unresolved & none & -0.74 & 0.001 & -1.72 & 0.009 & 62\% \\
  US loading rate & $+$ & -0.19 & 0.000 & +0.63 & 0.001 & 62\% \\
  EU discharge rate & $+$ & +0.84 & 0.001 & +0.46 & 0.000 & 62\% \\
  ballast returning to US & $+$ & +0.14 & 0.000 & +0.31 & 0.000 & 62\% \\
  laden voyage age & $-$ & -0.04 & 0.000 & -1.26 & 0.005 & 62\% \\
  US load-queue wait & $-$ & +1.01 & 0.001 & +0.55 & 0.001 & 38\% \\
  net export pressure & $+$ & -0.82 & 0.001 & -1.60 & 0.009 & 75\% \\
\bottomrule
\end{tabular}

\end{table}

\begin{table}[h]
\centering\footnotesize
\caption{Mechanism tests, full panel. The holdout coefficient for H2 changes sign
relative to the primary sample. In this sample 24\% of the holdout weeks lie after the
coverage seam, and the change of sign is not claimed as a finding.}
\input{tables/mechanismsfull}
\end{table}

\begin{figure}[h]
\centering
\includegraphics[width=\linewidth]{fig1_spread_and_stock_fullpanel}
\caption{Figure~\ref{fig:spread} extended to the data maximum. The
\coverageDropShare{} step in the at-sea stock at the start of 2026 is the coverage
seam and not a physical event.}
\end{figure}

\begin{figure}[h]
\centering
\includegraphics[width=0.9\linewidth]{fig3_fwl_scan_fullpanel}
\caption{Figure~\ref{fig:fwl} on the full panel.}
\end{figure}
